\documentclass[11pt]{article}
\usepackage{amsmath,amssymb,amsthm}
\usepackage[margin=25mm]{geometry}
\usepackage{hyperref}
\newtheorem{proposition}{Proposition}
\newcommand{\calP}{\mathcal P}
\newcommand{\calL}{\mathcal L}
\title{The original leading-coefficient line and its arithmetic relation image}
\author{Split Zero research programme}
\date{13 September 2026}
\begin{document}
\maketitle

This proof supplies the exact endpoint domain and metric of the multiplication
map in the paragraph following equation (13) of the delivered \emph{Arithmetic
Endpoint Bounds}. The arithmetic monic-norm theorem retains all of its original
constants and its domain $n\ge k\ge3$.

Fix a finite packet $h$, a positive integer $k$, and the original source map
\[
 \mathcal V_{h,k}P=P(D_1+\cdots+D_k)(F_h^{\otimes k}),\qquad
 D=-x\partial_x,\qquad \mathcal MF_h=(2\xi)/h.
\]
Its original Hermitian form, conjugate-linear in the first argument, is
\begin{equation}\tag{AL.1}
 \langle P,Q\rangle_{h,k}
 =\int_{\mathbb R}\overline{P(k/2+iu)}Q(k/2+iu)
    m_{h,k}(u)\,du,\quad
 m_{h,k}=w_h^{*k},\quad
 w_h(t)=\frac{|((2\xi)/h)(1/2+it)|^2}{2\pi}.
\end{equation}
The actual mass is $\int m_{h,k}=\mu_h^k>0$. The exponential upper bound
proved in the delivered arithmetic note makes every polynomial norm finite.
The density is positive almost everywhere; a nonzero polynomial has finitely
many zeros on the integration line. Thus (AL.1) is positive definite on every
finite polynomial space.

Let $\calP_a=\mathbb C[S]_{\le a}$ for $a\ge0$ and $\calP_a=0$ for $a<0$.
For $a\ge0$ put $\calL_a=\calP_a/\calP_{a-1}$ with its original quotient
norm. The leading-coefficient map is denoted $\operatorname{lc}_a$.
Let $p_a$ be the unique monic minimizer in (AL.1), and put
$\omega_a=\|p_a\|_{h,k}^2>0$. Its existence follows by projecting $S^a$
onto the finite-dimensional subspace $\calP_{a-1}$, and its uniqueness follows
from positivity. The defining orthogonality is
\begin{equation}\tag{AL.2}
 p_a\perp\calP_{a-1},\qquad
 \|z[p_a]\|_{\calL_a}^2=|z|^2\omega_a,\qquad z\in\mathbb C.
\end{equation}
Here $[p_a]=[S^a]$; this equality is in the degree quotient and does not identify
their polynomial representatives in the original source.

Fix the original monic arithmetic relation polynomial $\chi\in\mathbb C[S]$
of degree $q\ge1$ and the original remainder algebra
$E=\mathbb C[S]/(\chi)$. Write $\pi_\chi$ for polynomial reduction.

\begin{proposition}[Exact multiplication and leading-coefficient maps]
For $n\ge q-1$ the following sequence is exact:
\begin{equation}\tag{AL.3}
0\longrightarrow\calP_{n-q}\xrightarrow{\ \times\chi\ }
\calP_n\xrightarrow{\ \pi_\chi\ }E\longrightarrow0.
\end{equation}
For $n\ge q$ there is a commuting diagram of exact rows
\begin{equation}\tag{AL.4}
\begin{array}{ccccccccc}
0&\longrightarrow&\calP_{n-q-1}&\longrightarrow&\calP_{n-q}
 &\xrightarrow{\operatorname{lc}_{n-q}}&\mathbb C&\longrightarrow&0\\
 &&\downarrow\scriptstyle\times\chi&&\downarrow\scriptstyle\times\chi
 &&\downarrow\scriptstyle\mathrm{id}&&\\
0&\longrightarrow&\calP_{n-1}&\longrightarrow&\calP_n
 &\xrightarrow{\operatorname{lc}_n}&\mathbb C&\longrightarrow&0.
\end{array}
\end{equation}
It induces the coefficient-preserving isomorphism
\begin{equation}\tag{AL.5}
\overline m_\chi:\calL_{n-q}\xrightarrow{\sim}\calL_n,
\qquad [Q]\longmapsto[\chi Q].
\end{equation}
At $n=q-1$ the relation source in (AL.3) is zero and
$\pi_\chi:\calP_{q-1}\to E$ is an isomorphism. The degree map from that
relation source to $\calL_{q-1}$ is the zero map from $0$ to a
one-dimensional space.
\end{proposition}

\begin{proof}
Multiplication by a nonzero polynomial in the integral domain $\mathbb C[S]$
is injective. Its degree formula is
$\deg(\chi Q)=q+\deg Q$ for $Q\ne0$, so its image in $\calP_n$ consists
exactly of the multiples of $\chi$ in that space. Euclidean division by the
monic polynomial $\chi$ gives a unique remainder of degree below $q$.
That remainder belongs to $\calP_n$ precisely throughout the admitted range
$n\ge q-1$. This proves (AL.3), including its endpoint.

For $n\ge q$, both displayed leading-coefficient rows are exact by their
definitions, and the degree formula shows that the two multiplication maps
have the indicated targets. Since $\chi$ is monic,
\begin{equation}\tag{AL.6}
 \operatorname{lc}_n(\chi Q)=\operatorname{lc}_{n-q}(Q),
 \qquad Q\in\calP_{n-q}.
\end{equation}
This proves commutativity of the right square; the inclusion square commutes
because both composites are literal multiplication by $\chi$. Lower-degree
inputs map to $\calP_{n-1}$, so (AL.5) is well-defined. The two
leading-coefficient identifications with $\mathbb C$ take it to the identity
in (AL.4), proving that it is an isomorphism and specifying its inverse:
$z[S^n]\mapsto z[S^{n-q}]$.

At $n=q-1$, $\calP_{-1}=\calP_{-2}=0$. The exact endpoint diagram instead is
\begin{equation}\tag{AL.7}
\begin{array}{ccccccccc}
0&\longrightarrow&0&\longrightarrow&0&\longrightarrow&0
 &\longrightarrow&0\\
 &&\downarrow&&\downarrow&&\downarrow\scriptstyle 0&&\\
0&\longrightarrow&\calP_{q-2}&\longrightarrow&\calP_{q-1}
 &\xrightarrow{\operatorname{lc}_{q-1}}&\mathbb C&\longrightarrow&0.
\end{array}
\end{equation}
Every square commutes, and the degree target has dimension one even when
$q=1$, for which it is the constant line $\calP_0/0$. This proves the claimed
endpoint map and the exact domain restriction $n\ge q$ for (AL.5).
\end{proof}

\begin{proposition}[Exact source metric and arithmetic image]
For $n\ge q$, the norm, adjoint, and original polynomial representative of
(AL.5) are given by
\begin{equation}\tag{AL.8}
 \|\overline m_\chi\|^2=\frac{\omega_n}{\omega_{n-q}},\qquad
 \overline m_\chi^*[p_n]
 =\frac{\omega_n}{\omega_{n-q}}[p_{n-q}],\qquad
 \overline m_\chi[p_{n-q}]=[p_n].
\end{equation}
With the actual source correction
\begin{equation}\tag{AL.9}
 r_n=\chi p_{n-q}-p_n\in\calP_{n-1},
\end{equation}
one has the exact identities
\begin{equation}\tag{AL.10}
 \|\mathcal V_{h,k}(\chi p_{n-q})\|^2
   =\omega_n+\|\mathcal V_{h,k}r_n\|^2,\qquad
 \pi_\chi(r_n)=-\pi_\chi(p_n),\qquad
 \pi_\chi(\chi p_{n-q})=0.
\end{equation}
The unquotiented leading coefficient of every relation is transported by
(AL.6), and its arithmetic image is the original supported zero.
\end{proposition}

\begin{proof}
Both $\chi p_{n-q}$ and $p_n$ are monic of degree $n$. Their difference
therefore has degree at most $n-1$, proving (AL.9), and their classes in
$\calL_n$ agree. Formula (AL.2) gives
\[
 \frac{\|\overline m_\chi(z[p_{n-q}])\|^2}{\|z[p_{n-q}]\|^2}
 =\frac{|z|^2\omega_n}{|z|^2\omega_{n-q}}
 =\frac{\omega_n}{\omega_{n-q}}\quad(z\ne0).
\]
This proves the operator norm in (AL.8). The adjoint identity follows by
testing the defining equation against an arbitrary scalar multiple of
$[p_{n-q}]$ and using the positive real quotient weights in (AL.2).
Orthogonality $p_n\perp r_n$ gives the first identity of (AL.10) with the
literal original source norms. Reducing (AL.9) modulo $\chi$ gives its other
two identities.

To state the supported image with its complete type, for each original
admitted label $\lambda$ and the fixed admitted degree $n$, define the map
on $\{\tau\}\sqcup\{(\lambda,P):P\in\calP_n\}$ with target
$\{\tau\}\sqcup\{(\lambda,x):x\in E\}$ by
\[
 \widetilde\pi_\chi(\lambda,P)=(\lambda,\pi_\chi P),
 \qquad \widetilde\pi_\chi(\tau)=\tau.
\]
It sends a relation $(\lambda,\chi Q)$ to $(\lambda,0_E)$, while external
absence is fixed. The kernel inside this finite source fibre is exactly
$\chi\calP_{n-q}$, the image of multiplication in (AL.3). The same formula
on the full fibre $\mathbb C[S]$ has kernel $\chi\mathbb C[S]$;
its restriction along $\calP_n\hookrightarrow\mathbb C[S]$ is the finite
map just defined, because
$(\chi\mathbb C[S])\cap\calP_n=\chi\calP_{n-q}$ by the degree formula.
The original arithmetic map satisfies
\begin{equation}\tag{AL.11}
 J^{(k)}\mathcal V_{h,k}=\eta_{h,k}\pi_\chi,\qquad
 \eta_{h,k}[P]=\upsilon_h^{\otimes k}P(A_k)1,
 \quad \upsilon_h=j_h((2\xi)/h).
\end{equation}
Composing (AL.3) with this displayed original map sends each relation to
zero at that receiving fibre; every Taylor-unit coefficient stays in
$\eta_{h,k}$. Before that projection, (AL.10) records its full source norm
and the exact correction to the minimum-norm leading representative.
\end{proof}

For the empty arithmetic packet $\chi=1$, $q=0$, the algebra $E$ is zero.
Multiplication is the identity $\calP_n\to\calP_n$ and the induced line map
is the identity for every $n\ge0$, with norm one; (AL.8)--(AL.10) reduce to
$\omega_n/\omega_n=1$ and $r_n=0$. Its analytic source and actual mass
$\mu_h^k$ remain in (AL.1). This case uses no positive-rank quotient inverse.

The coefficient-preserving isomorphism, its original metric in (AL.8), the
representative correction in (AL.10), and the endpoint map (AL.7) give the
complete relation between the degree-line norm and its arithmetic relation
image. All identities hold at the original coordinate $S=k/2+iu$; no change
to the norm estimate or the source generator is required.
\end{document}
